\documentclass[12pt,a4paper,openright,oneside]{book} %Definisco carattere a 12 punti, foglio a4, inizio dei capitoli nella pagina a destra, 
%gestisce il pdf a 2 pagine
\usepackage[utf8]{inputenc} %Codifica utf-8 e mi permette di gestire i caratteri accentati.
\usepackage{disi-thesis} 
\usepackage{code-lstlistings}
\usepackage{notes}
\usepackage{shortcuts}
\usepackage{acronym}
\usepackage{newtxtext,newtxmath} %usa font times new roman
\usepackage{afterpage}

\school{\unibo}
\department{DIPARTIMENTO DI INFORMATICA -- SCIENZA E INGEGNERIA}
\programme{CORSO DI LAUREA IN INGEGNERIA E SCIENZE INFORMATICHE}
\scripttitle{Titolo TODO:}
\title{TODO}
\author{TODO}
\date{\today}
\subject{TODO}
\supervisor{TODO}
\cosupervisor{Dott. CoSupervisor 1}
\morecosupervisor{Dott. CoSupervisor 2}
\session{I}
\academicyear{TODO-TODO}

% Definition of acronyms
\acrodef{IoT}{Internet of Thing}
\acrodef{vm}[VM]{Virtual Machine}

\mainlinespacing{1.241} % line spacing in mainmatter, comment to default (1)

\begin{document}
\frontmatter\pagenumbering{arabic}\frontispiece
\tableofcontents  
\listoffigures
\listoftables
\lstlistoflistings
\afterpage{\null
           \thispagestyle{empty}
           \addtocounter{page}{-1}
           \newpage}
\begin{abstract}	
TODO
IoT (Internet of Things) è l'acronimo utilizzato per indicare la sfera tecnologica che ha alla base l'utilizzo di oggetti intelligenti. 
Per oggetti intelligenti si intendono quei dispositivi, macchine o attrezzature che si possono connettere a internet per trasmettere, ricevere dati o elaborarli.
All'interno di questa sfera, uno dei metodi di comunicazione più utilizzati è quello basato su un modello Publish/Subscribe che consente lo scambio di dati tra dispositivi della rete senza la necessità di connessioni dirette tra loro.
Per funzionare si utilizza il protocollo MQTT, uno dei protocolli più diffusi per lo scambio di messaggi, il quale si basa su una struttura con 3 tipologie di dispositivi connessi: Publisher, Subscriber e Broker.
I Publisher una volta ottenuti dei dati tramite sensori di vario genere, li pubblica, rendendoli disponibili nella rete. I dati per spostarsi all'interno della rete utilizzano i broker cioè vari punti/server di connessione, i quali permettono il passaggio di dati per poi infine arrivare ai Subscribers per la lettura e l'elaborazione.
Questa tesi riprende il lavoro effettuato da un collega nell'anno 21-22 approfondendo e sviluppando i casi di test con particolare attenzione a quelle che sono le caratteristiche che 
possono portare criticità al calcolo della soluzione ottima di percorso per il flusso dei dati.
La tesi del collega si basa sull'idea di costruire un sistema distribuito sfruttando le proprietà matematiche del rilassamento Lagrangiano, dove ogni nodo conosce solamente i propri vicini.

\end{abstract}
%\begin{dedication} % this is optional
%Optional. Max a few lines.
%\end{dedication} 
%----------------------------------------------------------------------------------------
\mainmatter
%----------------------------------------------------------------------------------------
\chapter{Introduction}
\label{chap:introduction}
%----------------------------------------------------------------------------------------

Write your intro here.
\sidenote{Add sidenotes in this way. They are named after the author of the thesis}

You can use acronyms that your defined previously,
such as \ac{IoT}.
%
If you use acronyms twice,
they will be written in full only once
(indeed, you can mention the \ac{IoT} now without it being fully explained).
%
In some cases, you may need a plural form of the acronym.
%
For instance,
that you are discussing \acp{vm},
you may need both \ac{vm} and \acp{vm}.

\paragraph{Structure of the Thesis}

\note{At the end, describe the structure of the paper}

\chapter{State of the art}

I suggest referencing stuff as follows: \cref{fig:random-image} or \Cref{fig:random-image}

\begin{figure}
    \centering
    \includegraphics[width=.8\linewidth]{figures/random-image.pdf}
    \caption{Some random image}
    \label{fig:random-image}
\end{figure}

\section{Some cool topic}

\chapter{Contribution}

You may also put some code snippet (which is NOT float by default), eg: \cref{lst:random-code}.

\lstinputlisting[float,language=Java,label={lst:random-code}]{listings/HelloWorld.java}

\section{Fancy formulas here}

%----------------------------------------------------------------------------------------
% BIBLIOGRAPHY
%----------------------------------------------------------------------------------------

\backmatter

\nocite{*} % Remove this as soon as you have the first citation

\bibliographystyle{alpha}
\bibliography{bibliography}

\begin{acknowledgements} % this is optional
Optional. Max 1 page.
\end{acknowledgements}

\end{document}
